\section{Introduction}
\label{sec:introduction}

Weak reversibility is a condition on the directed graph of a chemical
reaction network: each reaction has a directed return path. Two central
questions ask whether this condition alone prevents species extinction
under mass-action kinetics and whether every positive trajectory remains
bounded. We prove the joint boundedness and persistence conclusion for
arbitrary positive constant reaction rates, any number of connected
components in the
reaction graph, and every species dimension.

\subsection{The problem and main result}

A finite reaction network with $d\ge1$ species consists of a finite set of
\emph{complexes} $\mathcal C\subset\mathbb Z_{\ge0}^d$ and a set
$\mathcal R\subset\{(y,y'):y,y'\in\mathcal C,\ y\ne y'\}$ of directed
reactions. We write a reaction as $y\to y'$. The network is
\emph{weakly reversible} if, for each $y\to y'\in\mathcal R$, there is a
directed path from $y'$ to $y$. A connected component of the underlying
undirected reaction graph is called a \emph{linkage class}. There is no
restriction on their number. Empty reaction sets and isolated complexes
are allowed.

Assign a positive constant $\kappa_{y\to y'}$ to each reaction. The
mass-action equation is
\begin{equation}\label{eq:mass-action}
 \dot x=f(x):=
 \sum_{y\to y'\in\mathcal R}
 \kappa_{y\to y'}x^y(y'-y),
 \qquad x^y:=\prod_{i=1}^d x_i^{y_i}.
\end{equation}
Since the exponents are nonnegative integers, $f$ is a polynomial vector
field. A forward trajectory in $\R_{>0}^d$ is called \emph{persistent} if
\[
 \liminf_{t\to\infty}x_i(t)>0\qquad(1\le i\le d).
\]
The bounds in this definition may depend on the trajectory.

\begin{theorem}\label{thm:main}
Let $d\ge1$, let $\mathcal C\subset\mathbb Z_{\ge0}^d$ be finite, let $\mathcal R$ be
weakly reversible, and let all reaction rates be positive constants.
For every $x^0\in\R_{>0}^d$, the solution of \eqref{eq:mass-action} with
$x(0)=x^0$ exists for every $t\ge0$. There is a number
$\varepsilon\in(0,1)$, depending on the network, the rates, and $x^0$, such
that
\begin{equation}\label{eq:main-bounds}
 \varepsilon\le x_i(t)\le\varepsilon^{-1}
 \qquad(t\ge0,\ 1\le i\le d).
\end{equation}
\end{theorem}

In particular, every positive trajectory is bounded and persistent.
To relate the statement to conservation laws, put
\[
 S=\operatorname{span}_{\R}\{y'-y:y\to y'\in\mathcal R\},
 \qquad
 P=(x^0+S)\cap\R_{\ge0}^d.
\]
The subspace $S$ is the \emph{stoichiometric subspace}, and $P$ is the
closed \emph{stoichiometric compatibility class} of $x^0$. Since $f(x)$
takes values in $S$, integration of the equation shows that a positive
trajectory starting at $x^0$ stays in $P$.
Theorem~\ref{thm:main} applies whether or not $P$ is bounded. It therefore
includes the persistence question for bounded entire compatibility
classes. The estimate \eqref{eq:main-bounds} also shows directly that the
omega-limit set is nonempty and disjoint from the orthant boundary:
every sequence of times tending to infinity has a subsequence along which
the trajectory converges inside the compact positive box in
\eqref{eq:main-bounds}.
We make no assertion in this generality that the trajectory converges to an
equilibrium, or that the same $\varepsilon$ works for all positive initial
points.
In particular, the theorem does not assert uniform eventual bounds over
all positive initial points in a compatibility class.

The proof gives a geometric strengthening, stated precisely in
Proposition~\ref{prop:polytope}: each $x^0>0$ belongs to a compact convex
polytope $K\subset\R_{>0}^d$ such that every solution starting in $K$
exists globally and remains in $K$. The set $K$ need not be a whole
compatibility class. It is constructed in the ambient species space and
may meet several such classes. Its convexity also yields a short proof
of the known existence theorem of Boros~\cite{Boros2019}: every positive
stoichiometric class contains a positive equilibrium
(Corollary~\ref{cor:positive-equilibrium}).

\subsection{History and the role of the construction}

The classical structural theory begins with complex balance.
A system is \emph{complex-balanced} if it admits a positive state at
which the total incoming and outgoing reaction fluxes agree at every
complex. Horn and Jackson's equilibrium and stability
theory~\cite{HornJackson1972}, Horn's complex-balancing
criteria~\cite{Horn1972}, and Feinberg's deficiency-zero
analysis~\cite{Feinberg1972} provide the foundations. Weak reversibility
is a property of the graph, whereas complex balance also depends on the
rates. The global-attractor question for complex-balanced systems asks
whether every positive trajectory converges to the unique positive
equilibrium in its stoichiometric class. This convergence follows from
Theorem~\ref{thm:main} under the additional complex-balance assumption;
Remark~\ref{rem:complex-balanced-convergence} gives the argument.

Feinberg~\cite[Remark~6.1.E]{Feinberg1987} formulated the weaker
expectation that a positive trajectory of a weakly reversible mass-action system
cannot converge to a point where a species concentration is zero.
Persistence excludes approach to the boundary along any sequence of
times. Anderson~\cite[Section~1.1]{Anderson2011} separated the
\emph{boundedness conjecture}, that positive trajectories of weakly
reversible mass-action systems are bounded, from the
\emph{persistence conjecture} for bounded trajectories. Together these
imply persistence without an assumed upper bound.
Theorem~\ref{thm:main} establishes both conclusions for arbitrary
positive constant reaction rates.

August and Barahona~\cite[Theorem~6]{AugustBarahona2010} previously stated
the general boundedness and persistence conclusion. Their proof, Part~III
on page~45 before Equation~(11), infers monomial dominance from a maximal
sum of exponents and a sufficiently large product of concentrations.
That implication does not hold throughout the stated region. In
$2A\rightleftarrows B$ with both rates one, the concentrations
$(a,b)=(T,T^3)$ satisfy $ab\to\infty$ but $a^2<b$, and
\[
 \dot a+\dot b=b-a^2=T^3-T^2>0\qquad(T>1).
\]
This contradicts the asserted implication. It does not contradict their
theorem: $a+2b$ is conserved in this example. Our proof controls the
relevant exponent pairings separately and establishes the conclusion
independently.

Several earlier results give classwise permanence or permit
time-dependent rates. Here \emph{permanence} means that, for each fixed
compatibility class, there is a compact subset of the positive orthant
that eventually contains every positive trajectory in that class, with
entry time depending on the initial state.
Anderson~\cite{Anderson2011} proved boundedness for weakly reversible
systems with a single linkage class. Gopalkrishnan, Miller, and
Shiu~\cite{GopalkrishnanMillerShiu2014} proved permanence for the larger
class of strongly endotactic systems, which includes every weakly
reversible single-linkage system; Boros and
Hofbauer~\cite{BorosHofbauer2020} gave a later proof of the single-linkage
result. Craciun, Nazarov, and Pantea~\cite{CraciunNazarovPantea2013}
proved boundedness, persistence, and permanence for two-species weakly
reversible systems as part of their more general endotactic result.
Pantea~\cite{Pantea2012} proved persistence of bounded weakly reversible
trajectories when the stoichiometric subspace has dimension two, with
arbitrary species and linkage counts. Two species and two-dimensional
stoichiometry are different restrictions. These results also permit
rates that vary in time while remaining bounded above and away from
zero. Our theorem treats constant rates in all dimensions and with any
number of linkage classes. Its bounds may depend on the initial point,
so it does not assert the classwise uniformity of permanence.

The revised toric-differential-inclusion approach of
Craciun~\cite[Section~8, final remark, p.~81]{Craciun2026}, in version~3
dated September~23, 2026, gives persistence for bounded positive
trajectories of weakly reversible systems with uniformly positive and
bounded variable rates. Its complex-balanced application combines the
boundary construction with the Horn--Jackson Lyapunov bound. The general
weakly reversible statement assumes a trajectory upper bound.
Our construction supplies that bound together with persistence, through
a compact convex forward-invariant polytope containing each prescribed
positive initial point.

The geometric and analytical methods also have substantial antecedents.
Invariant polygons and quantitative monomial comparisons appear in
Craciun, Nazarov, and Pantea~\cite{CraciunNazarovPantea2013}, and
projected polyhedra in Pantea~\cite{Pantea2012}. Anderson's monomial
tiers and return-path domination~\cite{Anderson2011Global}, the
projection and multiscale arguments of Gopalkrishnan, Miller, and
Shiu~\cite{GopalkrishnanMillerShiu2013,GopalkrishnanMillerShiu2014},
and Craciun's ordered-complex embedding into toric differential
inclusions~\cite{Craciun2019} organize related comparisons. Our
finite affine-label construction and directed-cut estimate are proved
directly; none of these research results is an input to the proof.

\subsection{Proof overview}

The obstacle is to find a single finite family of linear inequalities
that controls both small and large concentrations. An individual reaction
need not point inward across a proposed face. Weak reversibility supplies
return paths, so the inward estimate will concern the combined flux across
cuts of the reaction graph. The normal of each active face must make
strict comparisons that agree with the monomial sizes. Near a tie, an
arbitrary approximation to a normal can reverse a comparison; the geometric
construction must respect these ties.

For $0<h<1$ and $p\in[-1,1]^d$, write
\[
 X_h(p)=(h^{p_1},\ldots,h^{p_d}).
\]
These points fill the box $B_h=[h,h^{-1}]^d$. We construct a concave
piecewise-affine function
\[
 F_h(x)=\min_{\lambda\in\Lambda}
       \bigl(r_\lambda\cdot x+c_\lambda(h)\bigr),
\]
where $\Lambda$ is a finite set indexing the affine functions, and the
slopes $r_\lambda\in[-1,1]^d$ are independent of $h$. A label
$\lambda$ is \emph{active} at $x$ if its affine function attains this
minimum; every minimizer is active at a tie.

Lemma~\ref{lem:affine} and Corollary~\ref{cor:uniform-affine} supply the
family for any positive tolerance $E:[-1,1]^d\to(0,\infty)$:
for all sufficiently small $h$, every label active at $X_h(p)$ satisfies
$\norm{p-r_\lambda}_\infty<E(r_\lambda)$, uniformly over the cube.
No continuity of $E$ is required. We can therefore make the tolerance
small near the hyperplanes where a coordinate or a complex comparison
vanishes, while keeping it positive on those hyperplanes. This forces
each nonzero slope comparison to have the same sign as the corresponding
exponent comparison, with a fixed positive gap. In particular, for some
$H>0$ independent of $h$,
\[
 r_\lambda\cdot y<r_\lambda\cdot z
 \quad\Longrightarrow\quad x^z\le h^H x^y
 \qquad(x=X_h(p),\ \lambda\text{ active}).
\]
The parameter order is essential: first the tolerance, then the fixed
finite slopes, then their positive separation $H$, and only then the
final value of $h$. A finite covering by nearby slopes alone would not
ensure that the slopes attaining the minimum have the required accuracy.

Order the complexes by an active slope. Across a cut between two
successive distinct levels, weak reversibility supplies an upward reaction
whenever there is a downward reaction. The monomial comparison makes one
upward rate dominate the sum of all downward rates once $h$ is small
enough. Summing over the cuts gives $r_\lambda\cdot f(x)\ge0$ for
every active label. Consequently $F_h$ is nondecreasing along any
trajectory segment in $B_h$. At a tie, its right derivative is the
minimum of the derivatives of all active branches, so controlling a
preferred active label would not suffice. Equal slope levels need not
have equal monomials; edges within one level contribute zero to the
normal derivative.

The coordinate comparisons force every label active at the prescribed
$x^0$ to have slope zero when $h$ is sufficiently small. Fixing such a
label gives a constant branch with value $M=F_h(x^0)$ and hence the
global bound $F_h\le M$. The maximum set
\[
 K=\{x\in\R^d:F_h(x)=M\}
\]
is an intersection of finitely many closed affine half-spaces and
contains $x^0$. A zero slope cannot be active on $\partial B_h$, so
$K$ misses that boundary. Convexity now places all of $K$ inside $B_h$:
a segment from $x^0$ to any point outside the box would meet its boundary.
Thus $K$ is a compact positive polytope. Along a solution starting in
$K$, nondecrease of $F_h$ and the ceiling $M$ force the solution to
remain in the maximum set. A first-exit argument and compact continuation
prove invariance and global existence, giving the trajectory bounds.
No active-label estimate outside $B_h$ is needed.

Section~\ref{sec:affine} proves the approximation lemma by induction on
dimension and finite gluing over the minimum exponent, then obtains its
uniform form. Section~\ref{sec:dynamics} imposes the sign constraints,
proves the cut and derivative estimates, constructs the invariant polytope,
and deduces the trajectory bounds, positive-equilibrium existence and
complex-balanced convergence.
